\section{Hạn chế}

\textbf{Về mô hình.} Phần lớn thực nghiệm so sánh các phương án phục hồi và
chưng cất mới được chạy ở quy mô sơ bộ trên CPU với tập COCO128, nhiều run dài bị
dừng sớm do giới hạn thời gian tính toán; kết luận định lượng cuối cùng phụ thuộc
vào pipeline đầy đủ trên tập PPE và phép đo trên thiết bị Jetson. Mới khảo sát một
tỷ lệ cắt tỉa chính (30\%); không gian tìm kiếm của SHO với ngân sách 6 cá thể
$\times$ 3 thế hệ và proxy 3 epoch còn nhỏ nên chưa bảo đảm tối ưu toàn cục.
Tập Construction-PPE chỉ có khoảng 1.400 ảnh, chưa phản ánh đủ điều kiện ánh sáng,
góc camera và khoảng cách ngoài thực tế.

\textbf{Về nền tảng.} Mỗi deployment nhắm tới một node; chưa có triển khai theo
nhóm, canary hay rollback tự động theo ngưỡng telemetry --- khôi phục hiện thực
bằng một lệnh update về model cũ. Chưa có kiểm tra checksum của model sau khi tải
ngoài việc ghim URL theo commit. Kết nối MQTT chưa bắt buộc TLS và chưa phân
quyền topic theo node trên broker. Xác thực dùng API key tĩnh; bảng người dùng,
API key băm và RBAC đã có trong schema nhưng chưa được dùng ở tầng ứng dụng.
Telemetry ngắn hạn trong Redis chỉ giữ 1 giờ. Mô-đun cảnh báo sớm mới dừng ở
mức thiết kế và đặc tả dữ liệu. Thu telemetry GPU dùng \texttt{nvidia-smi} nên chưa
có trên Jetson, cần bổ sung collector dựa trên \texttt{tegrastats}.

\textbf{Về AI Assistant.} Chất lượng phụ thuộc vào LLM thương mại và cần kết nối
Internet; cơ chế xác nhận thao tác ghi dựa trên prompt, chưa có luồng phê duyệt
tường minh và nhật ký kiểm toán riêng cho hành động của Assistant. Một kiểm thử
chuyển đổi kết quả công cụ ghi chưa đạt.
